% Options for packages loaded elsewhere
\PassOptionsToPackage{unicode}{hyperref}
\PassOptionsToPackage{hyphens}{url}
\PassOptionsToPackage{dvipsnames,svgnames,x11names}{xcolor}
%
\documentclass[
  10pt,
]{article}
\usepackage{amsmath,amssymb}
\usepackage{iftex}
\ifPDFTeX
  \usepackage[T1]{fontenc}
  \usepackage[utf8]{inputenc}
  \usepackage{textcomp} % provide euro and other symbols
\else % if luatex or xetex
  \usepackage{unicode-math} % this also loads fontspec
  \defaultfontfeatures{Scale=MatchLowercase}
  \defaultfontfeatures[\rmfamily]{Ligatures=TeX,Scale=1}
\fi
\usepackage{lmodern}
\ifPDFTeX\else
  % xetex/luatex font selection
  \setmainfont[]{DejaVu Serif}
  \setsansfont[]{DejaVu Sans}
  \setmonofont[]{DejaVu Sans Mono}
\fi
% Use upquote if available, for straight quotes in verbatim environments
\IfFileExists{upquote.sty}{\usepackage{upquote}}{}
\IfFileExists{microtype.sty}{% use microtype if available
  \usepackage[]{microtype}
  \UseMicrotypeSet[protrusion]{basicmath} % disable protrusion for tt fonts
}{}
\makeatletter
\@ifundefined{KOMAClassName}{% if non-KOMA class
  \IfFileExists{parskip.sty}{%
    \usepackage{parskip}
  }{% else
    \setlength{\parindent}{0pt}
    \setlength{\parskip}{6pt plus 2pt minus 1pt}}
}{% if KOMA class
  \KOMAoptions{parskip=half}}
\makeatother
\usepackage{xcolor}
\usepackage[margin=0.8in]{geometry}
\usepackage{longtable,booktabs,array}
\usepackage{calc} % for calculating minipage widths
% Correct order of tables after \paragraph or \subparagraph
\usepackage{etoolbox}
\makeatletter
\patchcmd\longtable{\par}{\if@noskipsec\mbox{}\fi\par}{}{}
\makeatother
% Allow footnotes in longtable head/foot
\IfFileExists{footnotehyper.sty}{\usepackage{footnotehyper}}{\usepackage{footnote}}
\makesavenoteenv{longtable}
\setlength{\emergencystretch}{3em} % prevent overfull lines
\providecommand{\tightlist}{%
  \setlength{\itemsep}{0pt}\setlength{\parskip}{0pt}}
\setcounter{secnumdepth}{5}
\usepackage{fvextra}
\DefineVerbatimEnvironment{Highlighting}{Verbatim}{breaklines,breakanywhere,commandchars=\\\{\},fontsize=\small}
\RecustomVerbatimEnvironment{verbatim}{Verbatim}{breaklines,breakanywhere,fontsize=\small}
\usepackage{microtype}
\usepackage{fancyhdr}
\pagestyle{fancy}
\fancyhf{}
\fancyhead[L]{\small Smaller Diophantine forms for Con(ZFC)}
\fancyhead[R]{\small Synthesis at 8c6eec9d19}
\fancyfoot[C]{\thepage}
\setlength{\headheight}{15pt}
\setlength{\emergencystretch}{4em}
\AtBeginDocument{\hypersetup{linkcolor=blue!45!black,urlcolor=blue!55!black}}
\makeatletter
\renewcommand*\l@section{\@dottedtocline{1}{0em}{2.4em}}
\renewcommand*\l@subsection{\@dottedtocline{2}{1.5em}{3.5em}}
\makeatother

\sloppy
\setlength{\tabcolsep}{3pt}
\ifLuaTeX
  \usepackage{selnolig}  % disable illegal ligatures
\fi
\usepackage{bookmark}
\IfFileExists{xurl.sty}{\usepackage{xurl}}{} % add URL line breaks if available
\urlstyle{same}
\hypersetup{
  pdftitle={Smaller Diophantine forms for ZFC consistency},
  pdfauthor={OpenAI assisted research for Vladimir Reshetnikov and ProveIt},
  colorlinks=true,
  linkcolor={Maroon},
  filecolor={Maroon},
  citecolor={Blue},
  urlcolor={Blue},
  pdfcreator={LaTeX via pandoc}}

\title{Smaller Diophantine forms for ZFC consistency}
\author{OpenAI assisted research for Vladimir Reshetnikov and ProveIt}
\date{23 October 2026}

\begin{document}
\maketitle

{
\hypersetup{linkcolor=}
\setcounter{tocdepth}{2}
\tableofcontents
}
\section{Results and interpretation}\label{results-and-interpretation}

This report combines the seven wave-4 submissions in ProveIt at commit
\path|8c6eec9d19ef80538c8e61ff6e0180040583446a|. It gives a small family
of explicit polynomials rather than claiming a single universal optimum.
Number of unknowns, permitted domain, total degree, arithmetic-circuit
size, and coefficient height measure different kinds of simplicity.

For every selected construction \(P\) below, with the fixed coefficients
and quantifier domain stated for that construction, the conclusion is

\[\mathrm{ZFC}\vdash
 \bigl(\neg\exists\mathbf x\;P(\mathbf x)=0\bigr)
 \longleftrightarrow \mathrm{Con}(\mathrm{ZFC}).\]

This does not decide either side. The new arguments are ordinary
mathematical proofs formalizable in ZFC; they are not new Lean or Rocq
certificates. The finite audits check polynomial identities, arithmetic
bounds, counts, fixed coefficient bytes, and transcription. They do not
replace the unbounded Pell/compiler arguments or prove that the machine
does not halt.

\begin{longtable}[]{@{}
  >{\raggedright\arraybackslash}p{(\columnwidth - 6\tabcolsep) * \real{0.2500}}
  >{\raggedleft\arraybackslash}p{(\columnwidth - 6\tabcolsep) * \real{0.2500}}
  >{\raggedright\arraybackslash}p{(\columnwidth - 6\tabcolsep) * \real{0.2500}}
  >{\raggedleft\arraybackslash}p{(\columnwidth - 6\tabcolsep) * \real{0.2500}}@{}}
\toprule\noalign{}
\begin{minipage}[b]{\linewidth}\raggedright
Domain and form
\end{minipage} & \begin{minipage}[b]{\linewidth}\raggedleft
Unknowns
\end{minipage} & \begin{minipage}[b]{\linewidth}\raggedright
Degree
\end{minipage} & \begin{minipage}[b]{\linewidth}\raggedleft
Binary operations
\end{minipage} \\
\midrule\noalign{}
\endhead
\bottomrule\noalign{}
\endlastfoot
Positive norm, cheapest evaluator & 12 & 1081 & \textbf{110} \\
Positive norm, recommended translation & 12 & \textbf{1041} &
\textbf{111} \\
Positive norm, lowest degree here & 12 & \textbf{1022} & 124 \\
Positive Jones compression & 9 & \(92324\cdot5^{32}+296\) &
\textbf{1554} \\
Nonnegative Jones compression & 9 & \(92324\cdot5^{32}+296\) & 1559 \\
Unrestricted integer compression & \textbf{14} &
\(813056\cdot5^{32}+2152\) & \textbf{1586} \\
Lower-degree integer compression & 15 & \(406528\cdot5^{32}+2528\) &
1594 \\
Unrestricted integer quartic & \textbf{50} & \textbf{4} & See residual
system \\
\end{longtable}

The new four-port radix is \(B=2^{1064921}\), compared with \(2^{5^9}\)
in the smallest wave-4 coefficient constructions. All four actual
coefficients occupy 532479 raw bytes. The new Jones index is
\(z=2^{4259933}\). The fourteen-integer polynomial eliminates the
enormous dense \(y\) coefficient altogether: its replacement is the
single explicit power of two \(2^{99184317514300346374515978685}\).

The most useful concise presentation is the twelve-positive-variable
form. The nine-variable construction minimizes the natural-domain
witness count among the explicit forms considered here but pays for it
with enormous degree and some enormous coefficient recipes. The
unrestricted-integer construction serves a different purpose: its
variables require no sign restrictions. The quartic is a low-degree
alternative, not a competitor on variable count.

\section{Comparison with the new
submissions}\label{comparison-with-the-new-submissions}

Packages 24--30 are mutually independent wave-4 submissions in the
pinned catalogue. Package 30 is our preceding construction. The
following comparisons use the packages' stated degree claims and
independently recounted live circuits where indicated in the
accompanying review ledger.

\begin{longtable}[]{@{}
  >{\raggedright\arraybackslash}p{(\columnwidth - 4\tabcolsep) * \real{0.3333}}
  >{\raggedright\arraybackslash}p{(\columnwidth - 4\tabcolsep) * \real{0.3333}}
  >{\raggedright\arraybackslash}p{(\columnwidth - 4\tabcolsep) * \real{0.3333}}@{}}
\toprule\noalign{}
\begin{minipage}[b]{\linewidth}\raggedright
Package
\end{minipage} & \begin{minipage}[b]{\linewidth}\raggedright
Main contribution
\end{minipage} & \begin{minipage}[b]{\linewidth}\raggedright
Selected previous figures
\end{minipage} \\
\midrule\noalign{}
\endhead
\bottomrule\noalign{}
\endlastfoot
24 & Tighter separating weights & 12 positive: degree 1069, 123
operations \\
25 & Mixed-radix layout and external population padding & Coefficient
radix \(2^{5^9}\) \\
26 & Ratio separation and parity copy tracks & 24 integers: degree 2432,
163 operations \\
27 & Multivariate Horner evaluation and exact coefficient aliases & Nine
positive: 1537 operations \\
28 & Signed-divisor finalizer & Nine nonnegative: degree
\(93264\cdot5^{33}+288\) \\
29 & Coprime Pell guards and intersective divisibility & 17 integers,
20607 operations \\
30 & Quartic register elimination and scalar norm evaluation & 12
positive: 111 operations; positive quartic in 31 variables \\
\end{longtable}

The new combinations are structural, not mere relabellings. They reuse
package 27's exact coefficient-expression aliases after reducing the
Jones parent to 31 variables; improve the signed majorants in package
28; replace the positivity mechanism used by package 29; and translate
the carrier in package 30's inexpensive norm. The fixed compiler tuple
is rebuilt coherently from the smaller-layout ideas, not assembled by
mixing ports from different submissions.

No construction here is claimed to beat the inherited 67-operation
single polynomial or 59-operation system on operation count alone. Those
remain useful alternatives. A system whose equality comparisons are free
is not directly comparable to a single-polynomial residual circuit.

\section{Conventions and proof
dependencies}\label{conventions-and-proof-dependencies}

We write \(\mathbb N_0=\{0,1,2,\ldots\}\) and
\(\mathbb Z_{>0}=\{1,2,\ldots\}\). All quantified variables are listed;
every other named expression is either an abbreviation or a fixed
integer coefficient. An arithmetic circuit uses only binary addition,
subtraction, and multiplication. Multiplication by a fixed numeral
counts as a gate. Shared subexpressions count once. Construction of
fixed coefficients is not included in the evaluation-gate count, and is
documented separately.

The common proof chain is:

\begin{enumerate}
\def\labelenumi{\arabic{enumi}.}
\tightlist
\item
  The fixed blank-tape run of K215 halts exactly when the specified
  syntactic ZFC proof system has a finite refutation.
\item
  The finite native tile language and its no-carry compiler describe
  exactly those halting runs.
\item
  The Pell arithmetic interface represents the compiler conditions.
\item
  The new polynomial transformations preserve existence of solutions in
  the explicitly stated domains.
\end{enumerate}

The first three steps are inherited mathematical dependencies, with
pinned source proofs supplied or identified in the archive. The fourth
step and the new compiler-parameter adapter are proved here. The
ZFC-provability claim is obtained by formalizing the arithmetic
arguments within ZFC, not by invoking Lean's or Rocq's soundness inside
ZFC.

Every arithmetic instance uses the same four ports \(\beta=B-1\),
\(K_0\), \(M_C\), and \(M_F\), where \(M_F\) is the shifted source mask,
not the native high mask. Whole-tuple bindings are mandatory. A smaller
radix does not reduce the Jones parent-variable count \(\nu\); that
count is 31 because of the separate quartic reduction.

\section{How the smaller fixed coefficients are
constructed}\label{how-the-smaller-fixed-coefficients-are-constructed}

\subsection{The theorem and its exact
scope}\label{the-theorem-and-its-exact-scope}

This construction combines three independently supplied ideas from the
seven wave-4 reports pinned at ProveIt commit
\path|8c6eec9d19ef80538c8e61ff6e0180040583446a|: package 25's mixed
clause fields and external zero tests, package 26's parity-separated
copy tracks, and package 24's Boolean controller for an arbitrary fixed
odd factor. The 943-tile language and its blank-machine interpretation
are unchanged.

The result is a complete four-port tuple with

\[
 b=121,\quad L=8801,\quad d=bL=1064921,\quad
 \rho=2^{121},\quad B=2^{1064921}.
\]

It satisfies the ordinary inputless native/Pell compiler interface,
including canonical square scale, odd strong multiplier and odd cubic
scale. Therefore it can instantiate the ordinary quartic31, the positive
twelve-variable reductions, the signed integer reductions, and then the
quartic-dependent Jones index construction. The last step must
regenerate its coefficients and index bound from the new tuple: its old
numerical Jones constants must not be retained.

This is a mathematical adapter to the inherited native/Pell and K215
calibration proofs, not a new Lean or Rocq certificate. No proof of
Con(ZFC) and no global minimum of coefficient height is asserted.

\subsection{Complete bounded recipe}\label{complete-bounded-recipe}

Take exactly the ordered943 tiles in \path|wang943_tiles.json|. They
have ordered fields (west,north,east,south). Start is selector0,
Forbidden is selector1. \path|build_coefficients.py| reconstructs all
943 tiles from all 430 literal K215 actions and checks equality of the
entire list, not just its length. The source machine SHA-256 is
\path|fdbcef0825886035c3091c8537a69535bc96f9e265f0ae45863543d40ae70090|.

Let \[
 e_t(j)=o_t+2j,\quad (o_0,o_1,o_2,o_3)=(1020,1000,980,960),
 \qquad W_j=2^{10+2j}.
\] For0\textless=j\textless10 these are the copy exponents of the four
edge labels. Define \[
 E=\{0,\ldots,942\}\cup\{e_t(j):0\le t<4,0\le j<10\}
 \cup\{1063,1064,1095,1096,2195,3296\}.
\] Thus \textbar E\textbar=989 and max E=3296. Each edge label x is
encoded by its ten-bit occupied word x+512. If tile i has labels x\_it,
let \[
 c_i=1+\sum_{t=0}^3\sum_{j=0}^9\operatorname{bit}_j(x_{it}+512)W_{10t+j}
       +\sum_{j=40}^{43}W_j.
\] The copy coefficient c\_(e\_t(j)) is W\_(10t+j). The coefficients at
anchors 1095,1096,2195,3296 are W40,W41,W42,W43 respectively; both dummy
coefficients c1063,c1064 are zero. Finally set \[
 \mu=1022+\sum_{j=0}^{43}W_j,
 \quad Z=(2^{119}-1)\sum_{e=8792}^{8798}\rho^e
          +(2^{27}-1)\rho^{8799}.
\] The complete native constants are \[
\begin{split}
 D_C={}&\rho^{40}+\rho^{41}+\rho^{1100}+\rho^{2200}
       +\sum_{e\in E}c_e(\rho^{4393-e}+\rho^{5495-e}),\\
 D_R={}&\rho,\\
 M_C={}&B-1-\sum_{e\in E,\ e\ne1}\rho^e-2\rho^{1063},\\
 M_{F,n}={}&4+\mu(\rho^{4393}+\rho^{5495})
       +\sum_{j=0}^9(\rho^{1000+2j}+\rho^{1021+2j})
       +\rho^{2195}+\rho^{3296}+Z.
\end{split}
\] The four polynomial ports, always replaced together, are \[
 \boxed{\beta=B-1,\qquad K_0=D_C+BD_R,\qquad M_C,
              \qquad M_F=M_{F,n}+B-1.}
\] MF is the shifted source high mask, not M\_(F,n). These are fixed
finite integer definitions; no coefficient is chosen by waiting for a
halt.

\subsection{Mixed clause decoding precedes payload carry
analysis}\label{mixed-clause-decoding-precedes-payload-carry-analysis}

Before decoding, n, the number of selected tile digits, lies
in{[}0,943{]}. Every nonconstant clause contribution is divisible by
1024. Consequently the low ten bits of the clause sum are exactly n,
even if there are carries among higher fields. Disjointness with the low
mask1022 forces n=0 or n=1. Only after that conclusion do we analyze the
44 payload fields. Each radix4 field is the sum of a Boolean
tile/occupancy bit and one Boolean copy/anchor bit, hence is 0,1,or2. No
payload carry remains, and masking its low bit permits0 or2. If n=0 all
copies/anchors vanish; if n=1 they equal the selected tile's occupied
labels and its four1anchors. This establishes exactly the same local
decoder as the old uniform-radix clause. It is not circularly assuming
all payload carries absent before extracting n.

\subsection{Exact geometry, arbitrary rotation and parity
separation}\label{exact-geometry-arbitrary-rotation-and-parity-separation}

For each j=0,\ldots,9 the complete convolution identities are:

\begin{itemize}
\tightlist
\item
  At1000+2j, DC sees960+2j through shift40; the temporal native word
  sees1000+2j; DR sees nothing.
\item
  At1021+2j, DC sees980+2j through shift41; DR sees1020+2j through
  shift1; the temporal word sees nothing.
\item
  At2195, DC sees1095 through shift1100, while the temporal word
  sees2195.
\item
  At3296, DC sees1096 through shift2200, while the temporal word
  sees3296.
\end{itemize}

In both anchor cases DR is absent. At each clause target4393,5495 the
coefficient on every basis digit e is exactly c\_e and DR is absent.
These are identities of the full DC coefficient polynomial, including
both reflected clauses, not bare-shift approximations. All69,230
component identities are checked exactly. Linearity makes them universal
in the native digits, including the harmless dummy values.

The parity tracks explain the short shifts: the temporal targets are
even, spatial targets odd, copy positions even, and other potential
predecessors are outside the native support. The reflected contributions
at the short clause locations are justified by the exact identities, not
by the inapplicable old inequality4393-max(E)\textgreater max(E).

The same exact finite support also gives \[
 (4393-E)\bmod L\;\cap\;(5495-E)\bmod L=\varnothing,
\] and \[
 (2195-E)\bmod L\;\cap\;(3296-E)\bmod L=\{0\}.
\] For the first, the band difference1102 is absent from E-E. For the
second, the only support pair of difference1101 is(2195,3296): the small
part has span1096, and the missing candidates1094 and 2194 prevent other
pairs. There are no additional circular cases because
L=8801\textgreater2max E. The upper-dummy spill position1064 is already
in E.

The inherited carry-free rotation argument supplies a clean clause for
any binary shift. Nonzero binary remainders1,\ldots,119 give only even
rotated digits; remainder120 permits at most the one upper-dummy odd
spill. Neither can supply the two required odd anchors in the marked
occupied cell. The binary remainder is therefore zero, and the anchor
intersection then forces a whole-cell rotation. The edge-bit equations
recover the exact Wang matches. The one-cell spatial shift propagates
occupancy from Start. No divisibility or coprimality restriction is
imposed on the temporal shift of an arbitrary purported solution.

\subsection{Uniform carries, positive slack and external zero
population}\label{uniform-carries-positive-slack-and-external-zero-population}

Let \[
 \mathcal M=991\left(2\sum_{e\in E}c_e+6\right).
\] Its bitlength is 118, while4M+8 has bitlength120. The clause mask has
bitlength97. The added padding is a mask, not an extra convolution term.
The explicit stronger estimates hold: \[
 \max(4\mathcal M+8,2\mu+4,2(2^{119}-1)+4,16)<2^{121},
\] \[
 3\rho/4+\mathcal M\le\rho-2,\qquad
 \rho/2+1+\mathcal M\le\rho-2.
\] These bounds apply before decoding, including all native selector
choices and dummy values, and retain the compiler's strict canonical
slack. The maximum unshifted DC-product exponent is 5495+3296=8791; DR
and temporal products are lower. Hence every bit tested by Z, at
8792,\ldots,8799, is identically zero for all aligned native
assignments. All positions are below L. Adding these tests cannot harm
completeness, and discarding them cannot harm soundness. No convolution
carry enters the padding.

The clause mask has 53 bits. The two clauses, twenty edge tests, two
anchor tests and low test contribute129 bits; Z
\path|contributes7*119+27=860| bits. Therefore \[
 \operatorname{pc}(M_{F,n})=989,
 \quad\operatorname{pc}(M_C)=d-989,
 \quad M_C\equiv2\pmod4,
 \quad M_{F,n}\equiv4\pmod8.
\] Both masks lie strictly between0 and B-1. These statements are
checked on the actual materialized integers. The two masks are not
claimed disjoint from one another: they constrain different packed
objects.

\subsection{Explicit rank controller without any power-of-five
hypothesis}\label{explicit-rank-controller-without-any-power-of-five-hypothesis}

This is the new completeness issue: d is no longer a power of five. Use
\[
 d=11^2\cdot13\cdot677,\quad \lambda=2341,\quad k=2340=\lambda-1.
\] Trial division proves all displayed primes. Exact modular arithmetic
gives \[
 \operatorname{ord}_{121}(B)=10,\quad
 \operatorname{ord}_{13}(B)=12,\quad
 \operatorname{ord}_{677}(B)=52,
\] so ord\_d(B)=780 divides k, lambda=1 modulo 780, and \[
 B^k\bmod\lambda^2=2556373\ne1,
 \qquad v_\lambda(B^k-1)=1.
\] The prime-unit rows (p,B mod p,(B\^{}2-1) mod p,H0 mod p), with
H0=1+B DC+B\^{}2(DR+1), are \[
 (11,2,3,7),\ (13,6,9,1),\ (677,611,293,53),\
 (2341,2100,1896,594).
\] Thus every needed factor is a unit at each prime of d lambda.

The unchanged tile language permits every sufficiently large spatial
width and temporal height independently: extra untouched blank margins
increase the width, and padded halting rows increase the height. This is
stronger than choosing powers of five in older convenient witnesses. The
precise pinned source is package 30/compiler/WANG943\_COEFFICIENTS.md,
Section 2.4, copied here as sources/WANG943\_LANGUAGE\_AND\_COMPILER.md;
it states arbitrary sufficiently large independent integer
widths/heights, not just powers of five. Choose width w\_c and height
h\_c to be sufficiently large powers of lambda, with
h\_c\textgreater lambda, N=w\_c h\_c=lambda\^{}n\textgreater d
lambda\^{}2, and q=B\^{}N. The recorded orders and Fermat's theorem give
\[
 q\equiv B^{w_c}\equiv B\pmod{d\lambda}.
\] Changing the free lower dummy in cell j changes the actual packed
rank by \[
 -\Gamma B^j,\quad
 \Gamma=\rho^{1063}(q^2-1)
         [1+q(D_C+BD_R+B^{w_c})].
\] The prime-unit table proves \path|gcd(Gamma,dN)=1,| since dN has no
other primes. Checking only a convenient five-adic residue would not
suffice here.

\subsubsection{Boolean coverage lemma}\label{boolean-coverage-lemma}

Put T=N/lambda and C0=d lambda. Every residue modulo C0 T=dN is the sum
of a Boolean subset of \[
 \{B^{2+kj}:0\le j<T\}\cup\{B^3\}.
\] Indeed, the non-Wieferich certificate and binomial lifting give exact
order T for B\^{}k modulo lambda\^{}n and order1 modulo d.~Its first T
powers are consequently exactly the T distinct residues1+C0 z modulo C0
T. Given target r, normalize by B\^{}-2. Choose epsilon in\{0,1\} such
that r B\^{}-2-epsilon B is nonzero modulo lambda. Let a be its least
residue modulo C0. Then1\textless=a\textless C0\textless T and
gcd(a,T)=1. Write the remaining target as a+C0 u and solve \[
 az_0+a(a-1)/2\equiv u\pmod T.
\] Select the a distinct orbit entries with additive coordinates
z0,z0+1,\ldots,z0+a-1 modulo T; include B\^{}3 if epsilon=1. Their sum
is the required residue. The extra cell3 is outside the orbit cells2+k j
because k\textgreater1. No cell is selected twice.

Every selected position is safe: the maximum orbit index is
N-N/lambda-lambda+3; adding1 or w\_c stays below N because
w\_c\textless N/lambda, and all positions are at least2. Thus none of
the required spatial/temporal translates wraps. Given the rank R0 with
these dummies zero, use the target(R0-d w\_c)Gamma\^{}-1. The new actual
rank satisfies R congruent d w\_c modulo dN, precisely the temporal
shift required by X=2\^{}R. Dummy changes preserve all decoded tiles and
mask tests, including the zero padding and the unique Start origin. The
uniform slack inequalities remain valid. This proves the new
completeness controller, not just its finite numerical prerequisites.

\subsection{Canonical parity promises and signed-family
transfer}\label{canonical-parity-promises-and-signed-family-transfer}

All d,w\_c,h\_c,N are odd. Write t=dN, so q=2\^{}t. The canonical unique
Start makes C congruent1 modulo 4. The native origin has no other
zero-exponent contribution; the temporal neighbor is not Start itself.
With \[
 S=(C-1)+qF,\quad M=q^2,\quad
 T_m=(M_CJ+1)+q(M_{F,n}J-1),
\] the native acceptance masks give S AND T\_m=0 and
0\textless=S,T\_m\textless M. Replicated blocks do not overlap. Since
M\_C=2 mod 4, adding1 to M\_C J increases its population by 1; since
M\_(F,n)=4 mod 8, subtracting1 from M\_(F,n)J also increases its
population by 1. Thus pc(T\_m)=N pc(M\_C)+N \path|pc(M_(F,n))+2=t+2.|
The exact source packing identity is \[
 R=(M-S)(M-1)+T_m=(M-1-S)M+(S+T_m).
\] Its high block is the 2t-bit complement of S; the low disjoint sum
adds populations. Hence \path|pc(R)=2t+pc(T_m)=3t+2.| Also S=0 mod 4 and
T\_m=3 mod 4, so R=3 mod 4. These facts persist through rank tuning.

The inherited strict bound R\textgreater3q+1 in particular gives
R\textgreater t. Canonical X=2\^{}R therefore has scale X/q=2\^{}(R-t),
a positive integer square because R and t are odd. The canonical main
Pell base a+2 is even. Its coordinate c=psi\_(a+2)(R) is odd, and the
same strong rank m=cR gives the normalized multiplier
\path|psi_(a+2)(m)/c^2| odd and positive. The ordinary Delta-scaled
multiplier is also odd since Delta=(a+2)\^{}2-1 is odd. These statements
are the inherited Pell binomial/parity lemmas with new odd t, not new
five-power assumptions.

Finally the inherited authentic half-binomial kernel identity gives
\path|v2(Y)=pc(R)-2=3t.| Thus Y=s q\^{}3 has positive odd s. This
supplies the extra completeness promise needed by the signed odd-scale
families. Soundness uses their unchanged weak transport bounds, Pell
rank recovery and signed root restoration; the four new ports satisfy
the entire native interface proved above. A coefficient transfer based
only on these congruences, without the rank controller and exact
geometry, would be insufficient.

\subsection{From the compiler interface to the requested logical
equivalence}\label{from-the-compiler-interface-to-the-requested-logical-equivalence}

The unchanged marked periodic tile language has a marked periodic tiling
if and only if the literal blank K215 computation halts. Soundness
traces its unique initialized head between reset rows; completeness pads
a finite halting computation with blank margins and halt rows. All these
finite arguments, the elementary Pell compiler, the controller above,
and the later polynomial eliminations are formalizable in ZFC. They
yield existence of a solution to any compatible fully instantiated
skeleton if and only if K215 halts, which the inherited finitary
calibration identifies with existence of a ZFC refutation. Taking
negations gives in ZFC the required equivalence between no solution and
Con(ZFC). No appeal to the soundness of Lean, Rocq, or ZFC itself is
used to assert that equivalence.

The ordinary quartic31 is an exact positive-zero-set elimination from
the ordinary source interface. Therefore its transfer follows from this
full interface, not merely from a match between normalized polynomial
ports. Jones's parameter nu remains31 for that quartic, regardless of
the much smaller native radix exponent d.~Changing d does not justify
changing nu.

\subsection{Size and verification
boundary}\label{size-and-verification-boundary}

The exact raw minimal unsigned little-endian payload sizes are:

\begin{longtable}[]{@{}lrr@{}}
\toprule\noalign{}
Port & Bitlength & Bytes \\
\midrule\noalign{}
\endhead
\bottomrule\noalign{}
\endlastfoot
beta & 1064921 & 133116 \\
K & 1065043 & 133131 \\
MC & 1064921 & 133116 \\
MF\_source & 1064922 & 133116 \\
\end{longtable}

The total is 532479 raw bytes, versus 977220 for packages 25and27; d is
45.476\% smaller than 1953125. Deterministic gzip storage in the
supplied build is 10466 bytes. Compression is not a coefficient-height
measure. The binding manifest records every raw/gzip SHA-256. Both raw
and compressed actual coefficients are provided; none of these four
ports remains a recipe only.

The builder and an independent checker verify all finite hypotheses,
including exact machine-to-tile equality, basis identities, native mask
populations, actual decoded constants, prime/order/lifting certificates,
and strict bounds. These finite audits supplement but do not replace the
universal mathematical proof or the inherited native/Pell calibration.

\section{A lower degree twelve variable
norm}\label{a-lower-degree-twelve-variable-norm}

\subsection{Results}\label{results}

All six new circuits use twelve strictly positive integer unknowns and
the same four \textbf{fixed}, coherently bound compiler coefficients.
The single-equation degree/operation pairs are

\begin{longtable}[]{@{}
  >{\raggedright\arraybackslash}p{(\columnwidth - 8\tabcolsep) * \real{0.1579}}
  >{\raggedleft\arraybackslash}p{(\columnwidth - 8\tabcolsep) * \real{0.2105}}
  >{\raggedleft\arraybackslash}p{(\columnwidth - 8\tabcolsep) * \real{0.2105}}
  >{\raggedleft\arraybackslash}p{(\columnwidth - 8\tabcolsep) * \real{0.2105}}
  >{\raggedleft\arraybackslash}p{(\columnwidth - 8\tabcolsep) * \real{0.2105}}@{}}
\toprule\noalign{}
\begin{minipage}[b]{\linewidth}\raggedright
Carrier correction
\end{minipage} & \begin{minipage}[b]{\linewidth}\raggedleft
Degree
\end{minipage} & \begin{minipage}[b]{\linewidth}\raggedleft
Binary operations
\end{minipage} & \begin{minipage}[b]{\linewidth}\raggedleft
System degree
\end{minipage} & \begin{minipage}[b]{\linewidth}\raggedleft
System operations
\end{minipage} \\
\midrule\noalign{}
\endhead
\bottomrule\noalign{}
\endlastfoot
no correction & 1081 & 110 & 1072 & 106 \\
b & \textbf{1041} & \textbf{111} & 1032 & 107 \\
\path|b+M(ac−ell)| & 1033 & 114 & 1024 & 110 \\
\path|b+M((a+2)c−ell)| & 1025 & 116 & 1016 & 112 \\
even-M, half-corrected & 1023 & 119 & 1014 & 115 \\
half-corrected plus auxiliary approximation & \textbf{1022} &
\textbf{124} & 1013 & 120 \\
\end{longtable}

All system entries have three equations. All operation counts charge
every binary addition, subtraction and multiplication, including by a
numeral, while sharing intermediate expressions. All degrees are exact,
not merely syntactic bounds. The first row is the previous package 30
polynomial with a cheaper evaluator. The other rows are different
polynomials whose zero sets are equisolvable, not claimed identical.

The principal recommendation is the 111-operation degree-1041 form: it
preserves the operation count of package 30's degree-1081 form and
improves package 24's degree-1069 form at 123 operations. The other rows
describe genuine further degree/cost tradeoffs, not a global minimum.
Twelve is not a minimum number of variables: the separate Jones
nine-variable reductions have much larger degrees.

\subsection{Fixed compiler interface}\label{fixed-compiler-interface}

Use one complete tuple (beta,K0,MC,MF) satisfying the normalized
blank-K215 interface. In particular MF is the shifted high mask. No
mixing of ports from different tuples is permitted. The release's
binding supplies the selected four fixed integers and their finite
recipe. The mathematical theorem below is uniform in any complete tuple
satisfying that interface, including package 30's materialized Wang943
tuple.

The supplied variables are

\begin{verbatim}
J,F,alpha,z,h,i,T,s,w,eta,zeta,t > 0.
\end{verbatim}

Retain the following definitions from package 30:

\begin{verbatim}
q=beta J+1, C=q−F−alpha, X=wq, Y=sq³, E=XY,
k=eta+zeta, c=kY+eta, a=E+Y, H=4a+3,
Delta=a²+H=(a+2)²−1,
ell=EkY, A=1+ell(ell+k), G=1+Delta c², b0=X+ac,
R=((q−1)(q−F)+alpha)(q²−1)+(MC+q MF)J,
Nk=k−hE−R, Nt=(K0+w)C+q−F−z(q−1),
I=ic², D=1+Delta I², Q=Delta(D−1)=(Delta I)²,
d=Q−1, v=cT, b=c+RD,
P=QDv², U=Qb², e=P+U−1, Baux=2Qvb,
M0=k(I²+T²).
\end{verbatim}

Every symbol other than the twelve supplied variables and four fixed
ports is a polynomial abbreviation, not an additional quantified
unknown.

For the first four rows put M=M0. For the last two put M=2M0. In every
case let

\begin{verbatim}
L=HM+1, W=b0+H(t+sigma).
\end{verbatim}

The corrections are respectively

\begin{verbatim}
sigma=0;
sigma=b;
sigma=b+M(ac−ell);
sigma=b+M((a+2)c−ell);
sigma=b+M0[2((a+2)c−ell)−k];
sigma=b+M0[2((a+2)c−ell)−k]+v(a+2)I.
\end{verbatim}

The fifth and sixth corrections can be written b+M delta, with

\begin{verbatim}
delta=(a+2)c−ell−k/2,
\end{verbatim}

but their displayed polynomial definitions contain no division: M=2M0
cancels the denominator. There are no variable exponents or rational
coefficients.

\subsection{Complete scalar
polynomial}\label{complete-scalar-polynomial}

Define computed expressions

\begin{verbatim}
n=(HM)² A, dn=d n,
HP=H²P, HU=H²U, He=HP+HU−H²,
j=4 HP HU,
Tr=2d(W²+L²G), Kr=d(W²−L²G), Kn=Kr²,
AA=2(dn+He), cc=(dn−He)²−j,
at=Tr−AA, bt=cc−Kn, co=16 dn j,
vr=at(at Tr+2bt)−co,
ur=bt²−at² Kn,
N=ur²+Tr ur vr+Kn vr².
\end{verbatim}

The promised polynomial and system are

\begin{verbatim}
Nk Nt(1+N)−1=0,
\end{verbatim}

or

\begin{verbatim}
Nk=1, Nt=1, N=0.
\end{verbatim}

This is a literal integer polynomial specified by a finite straight-line
circuit. The JSON files enumerate all gates and all supplied variables.
Squaring Kr rather than multiplying d² by (W²−L²G)² saves one binary
multiplication in every variant. This is an exact ring identity.

\subsection{Why the scalar expression is the
norm}\label{why-the-scalar-expression-is-the-norm}

For a formal argument x put

\begin{verbatim}
Cx=d(x²+M²A)−e,
F(x)=(Cx²+D Baux²−4d²x²M²A)²−4 Cx² D Baux².
\end{verbatim}

Then

\begin{verbatim}
N = H^16 F((W−L sqrt(G))/H) F((W+L sqrt(G))/H).
\end{verbatim}

The right side is a polynomial. More explicitly N is the product of the
eight factors

\begin{verbatim}
d(W−HM tau−L g)² − H²e + H² Baux f,
\end{verbatim}

where independently tau=±sqrt(A), g=±sqrt(G), f=±sqrt(D). Formal signs
here are a polynomial device and do not assert independent Galois
automorphisms.

To verify the scalar formula without field assumptions, use

\begin{verbatim}
(C²+j−K)²−4jC²=(C²−K−j)²−4Kj.
\end{verbatim}

After eliminating tau and f, the remaining polynomial in z has the shape

\begin{verbatim}
[(d z²)²−AA(d z²)+cc]²−co(d z²).
\end{verbatim}

For z=W±L sqrt(G), the two quantities x=d z² have sum Tr and product Kn.
Reduction modulo x²−Tr x+Kn gives ur+vr x. Its norm is ur²+Tr ur vr+Kn
vr². This proof works when any radicals specialize to zero, squares, or
coincident quadratic fields.

\subsection{Uniform soundness under positive carrier
corrections}\label{uniform-soundness-under-positive-carrier-corrections}

The inherited package 30 proof, included as
upstream/INHERITED\_POSITIVE12\_PROOF.md, establishes the following
facts from Nk=Nt=epsilon, epsilon in \{1,−1\}, before recovering
eliminated roots:

\begin{verbatim}
q≥16, X≥q, Y≥q³, E=XY,
3q<R<q⁴−2, k=R+epsilon+hE>E,
kY<c<2kY,
sqrt(G)−sqrt(A)>kY²,
2(v sqrt(D)+b)<M0 kY²< M0 sqrt(A).
\end{verbatim}

These are consequences of the entire authentic compiler tuple and
transport interface, not assertions for arbitrary ports. They imply for
any M≥M0 that every complex root x of F satisfies

\begin{verbatim}
|x|<M sqrt(G),
\end{verbatim}

and that F(0) is nonzero. The package 30 rational-recovery proof also
needs

\begin{verbatim}
d M² A>2Q(Dv²+b²),
\end{verbatim}

which follows from these same inequalities and is preserved by replacing
M0 with 2M0.

The new carrier leaves every one of these bounds unchanged. Its
correction sigma is positive whenever nonzero. Indeed b\textgreater0.
Also

\begin{verbatim}
ac−ell = kY²+a eta >0,
(a+2)c−ell >0,
(a+2)c−ell−k/2 > k(Y²−1/2)>0.
\end{verbatim}

Thus W\textgreater0 in every variant.

The identity Tr=2x0 and \path|Kn=x0²−4d²W²L²G,| with x0=d(W²+L²G), shows

\begin{verbatim}
N ≡ (ur+x0 vr)² (mod 4).
\end{verbatim}

At a solution of the product equation, 1+N cannot equal −1. The three
integer factors are units, so N=0 and Nk=Nt=epsilon with a common unit
sign. The three-equation system directly gives epsilon=1.

Choose a vanishing factor F((W−Lg)/H). If G is nonsquare, its two actual
quadratic conjugates would provide two roots of F with distance 2L
sqrt(G)/H, exceeding the common bound 2M sqrt(G) since L=HM+1. Thus g is
integral and x=(W−Lg)/H is rational. Since W\textgreater0 and
\textbar x\textbar\textless M sqrt(G), the selected g must be positive.

The inherited explicit rational-recovery formulas recover rational f and
tau with f²=D and tau²=A; both are integers. The rational auxiliary root
y=x−M tau obeys

\begin{verbatim}
Q(vf−b)²−(Q−1)y²=1.
\end{verbatim}

Because Q is an integer square, the rational Pell saturation lemma makes
y integral. Hence x is integral. The congruence W≡b0 (mod H), unchanged
by every new correction, gives g≡b0 (mod H), so gamma=(g−b0)/H is
integral. It is positive because g\textgreater(a+1)c and
c\textgreater X. The original normalized blank-machine interface,
including its signed quotient extension, now recovers the accepting K215
computation and epsilon=1 exactly as before.

All denominators and zero cases in this recovery are covered by the
inherited proof: Cx\textgreater0 for real x, Baux\textgreater0,
d\textgreater0, M\textgreater0, and F(0)≠0. No new rationality
assumption or sign change of an affine main-root expression is
introduced.

\subsection{Constructive completeness and positivity of the translated
variable}\label{constructive-completeness-and-positivity-of-the-translated-variable}

Take a canonical positive normalized blank-compiler witness, including
positive g,tau,y,gamma. Package30 uses

\begin{verbatim}
t_old=gamma+Mg+M tau+y.
\end{verbatim}

Set t=t\_old−sigma. This is an integer. Its defining identity gives
W=Lg+H(M tau+y), so N=0 and Nk=Nt=1. It remains essential to prove
t\textgreater0, rather than assuming a translation preserves the
positive domain.

For sigma=b, M tau\textgreater b follows from the bound above, and
t\textgreater0.

For sigma=b+M delta with delta=ac−ell or \path|delta=(a+2)c−ell| or its
half-corrected counterpart, one has delta\textless g. Indeed
ell/c\textgreater E/2\textgreater2, hence
(a+2)c−ell\textless ac\textless g; the other deltas are smaller.
Therefore

\begin{verbatim}
t=gamma+M(g−delta)+M tau+y−b>0.
\end{verbatim}

For the final auxiliary correction, use M=2M0 and

\begin{verbatim}
(a+2)I<2 sqrt(D).
\end{verbatim}

The latter follows from \path|D=1+((a+2)²−1)I²| and a+2≥2. Consequently

\begin{verbatim}
M tau >4(v sqrt(D)+b)>v(a+2)I+b.
\end{verbatim}

The same displayed formula for t, with the additional subtraction
v(a+2)I, is still strictly positive. This proves completeness without
enlarging a Pell index, imposing divisibility, or changing the canonical
coefficient calibration.

Thus all six constructions are equisolvable with blank halting of K215.
Composing with the inherited proof calibration and negating yields, in
ZFC, that their absence of solutions in the stated positive domain is
equivalent to Con(ZFC). This is a mathematical equivalence proof, not a
consistency proof or a new Lean/Rocq certificate.

\subsection{Exact degrees by the eight
factors}\label{exact-degrees-by-the-eight-factors}

All four fixed coefficients have degree zero. Scale every supplied
variable by a formal parameter Z and work in the Laurent-series field
over the rational function field of those variables and fixed
coefficients. The three positive-leading formal square roots exist
because their leading forms are squares. This is only a device for
reading the polynomial's homogeneous leading terms; no assumption about
numerical square roots is made.

The common degrees are

\begin{verbatim}
deg(E,H)=6, deg(Y)=4, deg(k)=1, deg(c)=5,
deg(ell)=11, deg(A,G)=22, deg(I)=11, deg(D)=34,
deg(Q,d)=46, deg(b)=38, deg(M0,M)=23.
\end{verbatim}

Write a bar over a symbol for its leading homogeneous form. Then

\begin{verbatim}
Hbar=4Ebar, cbar=kbar Ybar, ellbar=Ebar kbar Ybar,
M0bar=kbar Ibar², bbar=Rbar Dbar,
sqrt(D) has degree 17 and leading form Ebar Ibar.
\end{verbatim}

The positive-leading expansions satisfy

\begin{verbatim}
sqrt(G)−sqrt(A) = (ac−ell) + a term of degree 5,
sqrt(G)−sqrt(A)−(ac−ell) has leading form 2cbar,
sqrt(G)−sqrt(A)−((a+2)c−ell) has degree 1 and leading form −kbar/2,
((a+2)c−ell−k/2)−(sqrt(G)−sqrt(A))
    has degree−1 and leading form cbar/(2Ebar).
\end{verbatim}

The first difference sqrt(G)−sqrt(A) has degree 9 and leading form kbar
Ybar². These follow by expanding

\begin{verbatim}
sqrt(G)=sqrt((a+2)² c²−c²+1),
sqrt(A)=sqrt((ell+k/2)²+1−k²/4).
\end{verbatim}

For example \path|sqrt(G)=(a+2)c−c/[2(a+2)]+lower-degree| terms, while
\path|sqrt(A)=ell+k/2+terms| of degree at most−9. There is no vanishing
leading coefficient for authentic beta\textgreater0.

For the b-based variants write

\begin{verbatim}
Zsg,sa = W−Hb−HM sa sqrt(A)−L sg sqrt(G).
\end{verbatim}

The corresponding factor becomes

\begin{verbatim}
2dHb Zsg,sa + d Zsg,sa² − H²b² − H²QDv²
   + H² + sf H²Baux sqrt(D).
\end{verbatim}

Here the last radical term has degree 119. The four (sa,sg) pairs give
these degree lists, each repeated for the two f signs except in the last
row:

\begin{longtable}[]{@{}
  >{\raggedright\arraybackslash}p{(\columnwidth - 6\tabcolsep) * \real{0.2308}}
  >{\raggedright\arraybackslash}p{(\columnwidth - 6\tabcolsep) * \real{0.2308}}
  >{\raggedright\arraybackslash}p{(\columnwidth - 6\tabcolsep) * \real{0.2308}}
  >{\raggedleft\arraybackslash}p{(\columnwidth - 6\tabcolsep) * \real{0.3077}}@{}}
\toprule\noalign{}
\begin{minipage}[b]{\linewidth}\raggedright
Correction
\end{minipage} & \begin{minipage}[b]{\linewidth}\raggedright
Degrees of Zsg,sa
\end{minipage} & \begin{minipage}[b]{\linewidth}\raggedright
Degrees of the eight factors
\end{minipage} & \begin{minipage}[b]{\linewidth}\raggedleft
Sum
\end{minipage} \\
\midrule\noalign{}
\endhead
\bottomrule\noalign{}
\endlastfoot
b & 40, 40, 38, 38 & 130, 130, 130, 130, 128, 128, 128, 128 & 1032 \\
\path|b+M(ac−ell)| & 40, 40, 38, 34 & 130, 130, 130, 130, 128, 128, 124,
124 & 1024 \\
\path|b+M((a+2)c−ell)| & 40, 40, 38, 30 & 130, 130, 130, 130, 128, 128,
120, 120 & 1016 \\
half-corrected & 40, 40, 38, 28 & 130, 130, 130, 130, 128, 128, 119, 119
& 1014 \\
half-corrected plus auxiliary & 40, 40, 38, 29 & 130, 130, 130, 130,
128, 128, 119, 118 & 1013 \\
\end{longtable}

In the first three rows each leading factor is 2 dbar Hbar bbar times
its nonzero Z-leading form. In the half-corrected row the two special
leading terms are the opposite nonzero f terms of degree 119.

For the last row the special branch is \path|(sa,sg)=(−1,+1).| Put
f0=(a+2)I and \path|delta=(a+2)c−ell−k/2.| There

\begin{verbatim}
Zsg,sa=H v f0 + H M(delta+sqrt(A)−sqrt(G)) + b0+Ht−sqrt(G).
\end{verbatim}

Its degree 29 leading term is Hbar vbar fbar. At one f sign the two
degree 119 contributions add; at the other they cancel. The next
surviving degree 118 term is

\begin{verbatim}
2 dbar Hbar² bbar Mbar cbar/(2Ebar),
\end{verbatim}

which is nonzero. To check that no omitted term has degree 118 or
greater: f0−sqrt(D) has degree 5, so the correction from replacing f0 by
sqrt(D) contributes degree 107; b0+Ht−sqrt(G) has degree at most7; and
all remaining square terms have degree at most104. Thus the final factor
has exactly degree 118.

The product degree is the sum of the eight nonzero factor degrees in
this Laurent field, and the product is the integer polynomial N. This
proves all five exact norm degrees. For the unshifted baseline, the
inherited proof gives N\_top=Hbar\^{}16 Ubar\^{}8≠0 and degree 1072.
Finally Nk and Nt have exact degrees 7 and 2 with nonzero leading forms,
so the product polynomial has degree deg(N)+9. Every claim holds for the
actual fixed tuple, not merely for generic coefficient values.

\subsection{Verification and scope}\label{verification-and-scope}

The new emitter contains no received executable imports. Its
accompanying checker independently multiplies the eight radical factors
using an eight-dimensional formal algebra; it checks positive, negative,
zero, and degenerate inputs modulo several primes. Additional tests
check the unit congruence, all gates and ports live, the factorization
underlying the cheaper Kn evaluation, the exact degree claims by full
finite-field univariate polynomial arithmetic, and synthetic
radical-kernel witnesses. These finite diagnostics support the displayed
identities but do not replace the mathematical proofs or verify the
inherited compiler by finite testing.

Sources are pinned at commit
\path|8c6eec9d19ef80538c8e61ff6e0180040583446a|. Package24 motivates
deliberate leading-term cancellation; package 30 supplies the scalar
trace norm and the separating bound. The new idea is translating the
already quantified carrier by increasingly accurate polynomial
approximations to the auxiliary and radical expressions while proving
that the translated positive variable remains positive.

\section{A quartic in fifty unrestricted
integers}\label{a-quartic-in-fifty-unrestricted-integers}

The previous signed lift has 51 integer coordinates and 30 quadratic
residuals. Its only use of a packing register \(G\) is

\[G=q(q-F)-C,\qquad
k-1-hE=LG+(M_C+qM_F)J,\]

where \(q=\beta J+1\), \(C=q-F-\alpha\), and the retained equations
include \(L=q^2-1\) and \(N=q(L+1)\). Delete \(G\) and replace the
packing equation by

\[k-1-hE=(q-F)(N-q)-CL+(M_C+qM_F)J.\]

It is still quadratic. The polynomial identity

\[L[q(q-F)-C]-[(q-F)(N-q)-CL]
  =(q-F)[q(L+1)-N]\]

proves an exact bijection of the integer zero sets: forget \(G\) in one
direction, and restore \(G=q(q-F)-C\) in the other. Unlike the
positive-domain version, no positivity proof for this restored register
is needed.

The other 28 residuals, including all six positivity guards and the odd
scale, remain unchanged. Hence this reduction inherits exactly the
source's signed-domain soundness theorem. Taking the sum of the 29
squared residuals gives a quartic in 50 unrestricted integers. Its
coefficient of \(\tau^4\) is 1, so its degree is exactly four for every
admissible coefficient tuple. The archive contains all residuals, a
symbolic identity check, and a separate independent reconstruction
audit.

\section{Nine natural variables with a smaller signed
finalizer}\label{nine-natural-variables-with-a-smaller-signed-finalizer}

The local majorant MF below is unrelated to the fixed compiler mask
port.

\subsection{Result, scope, and credits}\label{result-scope-and-credits}

Set L=5\^{}32. There is an explicitly specified fixed polynomial with
nine nonnegative integer unknowns, exact total degree

\begin{verbatim}
92324 L + 296 = 2149585634469985961914062796,
\end{verbatim}

and a division-free circuit with 1,559 binary operations (692 additions,
22 subtractions, 845 multiplications). An equivalent polynomial over
nine strictly positive integer unknowns has 1,554 operations. Both have
80 nonliteral fixed leaves: the two index integers and 78 distinct
nonliteral coefficient recipes. The 402 original weighted coefficient
records remain the canonical coefficient specification.

These combine package 30's quartic31, package 27's multivariate Horner
and coefficient-recipe folding, and package 28's signed-divisor
finalizer. The new mathematical reductions are the signed quadratic
majorants below, which cause a provable leading cancellation and lower
the degree, and replacing u(2m−1) by um in the finalizer. Sources are
pinned at \path|8c6eec9d19ef80538c8e61ff6e0180040583446a|.

The claimed equivalence is a mathematical theorem formalizable in ZFC:
no solution in the stated domain if and only if Con(ZFC). It does not
prove consistency, decide existence, assert global minimality, or
provide a new proof-assistant certificate. It imports the explicit
quartic/Pell compiler interface, the coherent coefficient adapter, K215
calibration, and the constructive parameterized Jones reduction. The new
arithmetic steps are proved below without assuming any radicand is a
square.

\subsection{Fixed coefficient
instance}\label{fixed-coefficient-instance}

The parent is the 31-positive-variable, 22-quadratic system in
quartic31.json. Its sum of squared residuals after replacing each
positive coordinate by n\_i+1 is P31(x,n), with unused dummy input x. It
has degree four (coefficient of n\_tau\^{}4 is one) and P31(x, 0)≥1 (the
residual EX+E−aX is one at the all-ones tuple). The register-elimination
proof from package 30 is retained as QUARTIC31\_PROOF.md. All four
coefficient ports are replaced together by
../coefficients/payloads/COEFFICIENT\_BINDING.json:

\begin{verbatim}
beta=2^1064921−1, K=Kconstant, MC, and MF=MF_source.
\end{verbatim}

The fourth is the shifted mask, never the native high mask. The complete
coefficient adapter, not individual modular identities, establishes that
this tuple implements the ordinary positive compiler interface.

For each \path|alpha=(alpha_0,...,alpha_31)| of weight at most four, set
\path;c_alpha=(4−|alpha|)!; product\_j alpha\_j! times the coefficient
of x\^{}alpha\_0 n\_1\^{}alpha\_1 \ldots{} n\_31\^{}alpha\_31 in P31.
Set \path|expo(alpha)=sum_j| alpha\_j 5\^{}j. All 402 nonzero symbolic
coefficients in the four ports are listed in INDEX\_CONSTANTS.json. The
other coefficients are zero. There are \path|binomial(36,4)=58,905|
eligible alpha.

Let s=4,259,933 and z=2\^{}s. INDEX\_BOUND.json records a direct finite
bound: if c\_alpha=sum\_t a\_t kappa\^{}e\_t and the four port bit
lengths are b\_j, then \textbar c\_alpha\textbar{} ≤ sum\_t
\textbar a\_t\textbar{} 2\^{}(sum\_j b\_j e\_tj), and this last integer
plus one is strictly smaller than 2\^{}s for every record. Here the bit
lengths are (1,064,921; 1,065,043; 1,064,921; 1,064,922). Thus

\textbar c\_alpha\textbar+1 \textless{} z.

Define \path!e(T)=sum_{|alpha|≤4}! (z+c\_alpha)
\path|T^(L−expo(alpha)),| y=e(2z), and the proof-interface constant
\path|u_index=sum_{j=1}^31| (2z)\^{}(5\^{}j). All exponents are positive
integers, since \path|expo(alpha)≤4*5^31<L.| Base-5 uniqueness gives
distinct digit positions. Each digit z+c\_alpha lies strictly between
one and 2z. Consequently

2(2z)\^{}L ≤ y \textless{} (2z)\^{}(L+1).

The left inequality uses c\_0=24 P31(0)≥24. These facts are the explicit
Jones index hypotheses. The auxiliary u\_index does not occur in the
polynomial. z is materialized as a 532,492-byte little-endian integer;
All 78 distinct nonliteral weighted coefficient leaves are also
materialized as signed-magnitude gzip files: 22,496,535 raw magnitude
bytes compress to 2,427,559 bytes. Only y remains an explicit bounded
recipe rather than a materialized integer. The 402 canonical coefficient
recipes remain available for verification and aliases. No uncomputed
recursively enumerable index or unspecified integer is hidden in a
coefficient.

\subsection{Complete polynomial
description}\label{complete-polynomial-description}

The nine unknowns are epsilon,g0,h0,s0,w0,phi,i0,j0,n in
N=\{0,1,\ldots\}. Set \path|g=g0+1,h=h0+1,s=s0+1,w=w0+1,i=i0+1,j=j0+1,|
and

\begin{verbatim}
b=y+1+epsilon, B=b^5, q=B^L,
ell=sum_{k=1}^31 B^(5^k), lambda=sum_{k=0}^{4L−1}B^k,
N=q^16, Gcap=32 b B^(5^31).
\end{verbatim}

Define the packed polynomial

\begin{verbatim}
R=[g+e(B)q^3+ell q^5+
 (2(e(B)−z lambda)(1+B+g)^4+lambda B+lambda B q^4)q^7](N^2−N)
 +[q^3−b ell+ell+(B−2z)lambda q^3+(B−2)q^8](N^2−1).
\end{verbatim}

Continue with

\begin{verbatim}
U=N^2 w, Y=N^2 s, M=RY,
P=2M^2 U, K=R+1+h(P−1), A=M(U+1),
B1=2R+1, W1=bw, C=B1+W1+phi,
J=4A−5, Q=3W1 C−2(W1^2−1),
D=(A^2−1)C^2+1, E=i J D C^2,
F=(A^2−1)E^2+1, G=A+F(CD+1−A),
H=B1+jC, I=(G^2−1)H^2+1,
a0=(P^2−1)K^2+1, a1=DFI,
u=JF, v=QF+J(H−C),
m=K^2 Gcap−4(C−KY)^2 Gcap−8gK^2.
\end{verbatim}

Every symbol is a polynomial abbreviation, not an additional unknown.
The letter u here is the signed divisor and is distinct from u\_index.
Now put

\begin{verbatim}
r=W1+phi, sigma=Y(U+1),
MD=sigma(RC+r^2), TC=CH+r^2,
MF=AE+1,
MI=F(D TC+2H^2+A^2)−A^2+1,
V0=PK+1, V1=MD MF MI+1,
Mdiv=Q^2+(H−C)^2+1,
Omega=u m,
Theta=u(n+Mdiv)+v−Omega(2Mdiv+V0 V1).
\end{verbatim}

The requested polynomial is

\begin{verbatim}
F9=(Theta^2−Omega^2(a0+V0^2 a1))^2
       −4 Omega^4 V0^2 a0 a1.
\end{verbatim}

All powers have fixed nonnegative exponents. The JSON circuit expands
each power into binary multiplication gates. Its natural-domain filename
is nine31\_signed\_quadratic.json. For strictly positive unknowns
substitute each supplied variable by its value minus one; the optimized
translated circuit is nine31\_signed\_quadratic\_positive.json.

\subsection{The packed integer cannot
vanish}\label{the-packed-integer-cannot-vanish}

The index gives b≥2. Reduction modulo q gives

R ≡ (b−1)ell (mod q).

Indeed q divides N, and every term of the second bracket except (1−b)ell
is divisible by q. Moreover

0 \textless{} (b−1)ell ≤ (B−1)ell ≤ \path|(B−1)sum_{k=0}^{L−1}B^k| =q−1.

Hence R≠0 for every nonnegative assignment, before imposing F9=0. Since
U,Y≥1, sigma≥2 and A=R sigma imply \textbar A\textbar≥2. Also
P≥2R\^{}2≥2 and K≥R+P≥2R\^{}2+R≥1. Thus J≠0, D≥1,
AJ=A(4A−5)\textgreater0, AE=iAJDC\^{}2≥0, and F≥1. In particular u=JF≠0.

\subsection{Signed quadratic
majorants}\label{signed-quadratic-majorants}

Here r=bw+phi≥2, C=2R+1+r, and H=(j+1)C−r. The elementary identity

\path|2RC+r^2=4(R+(r+1)/4)^2+r^2−(r+1)^2/4|

is strictly positive when r≥2. It follows that
RC+r\^{}2\textgreater\textbar RC\textbar: if RC≥0 its excess is r\^{}2,
and if RC\textless0 the excess is 2RC+r\^{}2. Since the quantities are
integers, the excess is at least one. Therefore

\begin{verbatim}
MD=sigma(RC+r^2) ≥ |AC|+sigma ≥ |AC|+2 ≥ sqrt(D).
\end{verbatim}

In particular MD≥1. This statement permits either sign of R and C, and
C=0. No assumed sign of m or square condition was used.

Likewise

\path|2CH+r^2=2(j+1)C^2−2rC+r^2≥0|

by completing the square (j≥1). Thus TC=CH+r\^{}2≥\textbar CH\textbar.
The coefficient j+1 is why this bound remains valid for negative C and
for H=0. Since AE≥0,

\begin{verbatim}
MF=AE+1≥1,   F=(AE)^2−E^2+1≤MF^2.
\end{verbatim}

Using G=FDC+F−A(F−1), integer inequalities \textbar H\textbar≤H\^{}2 and
\textbar AH\textbar≤A\^{}2+H\^{}2 yield

\textbar GH\textbar{} ≤ FD TC+F \path|H^2+(F−1)(A^2+H^2).|

Adding \textbar H\textbar+1 and using \textbar H\textbar≤H\^{}2 proves

\begin{verbatim}
MI=FD TC+2F H^2+(F−1)A^2+1 ≥ |GH|+|H|+1.
\end{verbatim}

Consequently MI≥1 and \path!sqrt(|I|)≤sqrt(G^2H^2+H^2+1)≤|GH|+|H|+1≤MI.!
This includes negative I and both zero edge cases. Hence

V0,V1≥2, sqrt(\textbar a0\textbar)≤V0−1, sqrt(\textbar a1\textbar)≤V1−1.

For a0 use \path|1≤a0=(PK)^2−K^2+1≤(PK)^2.| For a1 multiply D≤MD\^{}2,
F≤MF\^{}2, \textbar I\textbar≤MI\^{}2. These inequalities are
unconditional on the natural outer domain.

\subsection{A zero-margin-safe signed-divisor
finalizer}\label{a-zero-margin-safe-signed-divisor-finalizer}

Lemma. Suppose a0,a1,u,v,m,M,V0,V1 are integers with u≠0, V0,V1≥2,
sqrt(\textbar at\textbar)≤Vt−1, M≥1, and \textbar v/u\textbar≤M−1. For
n≥0 set Omega=um and \path|Theta=u(n+M)+v−Omega(2M+V0V1).| Then the
polynomial in Section 7.3 has a nonnegative integer n making it zero if
and only if a0,a1 are integer squares, u divides v, and m\textgreater0.

Proof. If m=0, Omega=0 and \path|Theta/u=n+M+v/u≥n+1>0.| Thus the
polynomial is Theta\^{}4≠0. Otherwise Omega≠0. Choose complex square
roots r0\^{}2=a0, r1\^{}2=a1. The polynomial factors exactly as
product\_\{epsilon,delta in \{−1,1\}\}{[}Theta−Omega(epsilon r0+delta
V0r1){]}. A zero allows signs to be absorbed so that
\path|s=r0+V0r1=Theta/Omega| is rational. If s≠0 then 2s
\path|r0=s^2+a0−V0^2a1| makes both roots rational and hence integers
(their squares are integers). If s=0, \path!|r1|=|r0|/V0≤(V0−1)/V0<1!
forces the integer a1 to have absolute value below one, so a1=a0=0. This
avoids any independence assumption on quadratic fields and handles
negative radicands.

Let \path|Z=2M+V0V1+r0+V0r1.| The root bounds give Z≥2M+1\textgreater0.
The vanishing factor becomes n+M+v/u=mZ. Its left side is positive, so
m\textgreater0. Its right side is integral, so v/u is integral and u
divides v. Conversely choose integer roots and let n=mZ−M−v/u. Then n is
an integer and n≥2M+1−M−\textbar v/u\textbar≥2. Its factor vanishes.
QED.

For the Jones application, \path!|v/u|=|Q/J+(H−C)/F|!
\path!≤|Q|+|H−C|≤Q^2+(H−C)^2=Mdiv−1.! The lemma therefore applies. The
choice Omega=um saves one multiplication and one subtraction relative to
u(2m−1); the extra m=0 case above is essential.

\subsection{Equivalence to
consistency}\label{equivalence-to-consistency}

For each fixed eight-coordinate outer assignment, the last lemma
identifies existence of n with exactly the inherited four Jones ratio
conditions: a0 and a1 squares, JF dividing QF+J(H−C), and
m\textgreater0. The constructive parameterized Jones packing and ratio
theorems then identify existence of all nine coordinates with a zero of
P31(1,·). This uses nu=31, not the separate public existence theorem
normalized at nu=58. The explicit index in Section 7.2 satisfies the
former theorem's hypotheses. The quartic31 algebraic reduction and
complete coefficient adapter identify those zeros with a finite blank
K215 halting computation; the K215 calibration identifies this with a
finite ZFC refutation. All transformations concern finite arithmetic
objects and are provable in ZFC. Negating their existence equivalence
gives precisely

\begin{verbatim}
ZFC proves: [no N^9 solution to F9=0] iff Con(ZFC).
\end{verbatim}

The positive-domain translated circuit has the same conclusion by the
bijection between N and positive integers. Neither statement is valid
merely by reading the same nine-variable formula over unrestricted Z.

\subsection{Exact degree despite a top
cancellation}\label{exact-degree-despite-a-top-cancellation}

Degree counts the nine unknowns, with fixed coefficients of degree zero.
On the diagonal where every unknown is t, relevant degrees and leaders
are as follows (all leaders are independent of fixed coefficients):

R:235L, 1; A:555L+2,1; C:235L, 2; D:1580L+4,4; E:2605L+7,2\^{}6;
\path|F:6320L+18,2^12;| \path|G:8135L+22,2^15;| H:235L+1,2;
\path|I:16740L+46,2^32;| MD:790L+2,2; MF:3160L+9,2\^{}6;
\path|MI:8370L+23,2^16;| V0:1900L+7,4; \path|V1:12320L+34,2^23;|
\path|Omega:9096L+31,−2^23;| \path|Theta:23316L+72,2^48.|

The naive top terms in Theta\^{}2−Omega\^{}2 V0\^{}2 a1 cancel exactly.
To prove the true degree, use the following polynomial identities:

\begin{verbatim}
DeltaD=MD^2−D=2AC sigma r^2+sigma^2 r^4+C^2−1;
DeltaF=MF^2−F=E^2+2AE;
DeltaIRoot=MI−GH=FD r^2+2FH^2+(F−1)A^2+1−FH+A(F−1)H;
DeltaI=MI^2−I=DeltaIRoot(MI+GH)+H^2−1;
Delta=V1^2−a1=DeltaD MF^2 MI^2+D DeltaF MI^2+DF DeltaI+2MD MF MI+1.
\end{verbatim}

The leading term of DeltaD has degree 1110L+8 and coefficient4. The
first summand in Delta therefore strictly dominates, giving Delta degree
24170L+72 and \path|coefficient2^46.| Let \path|E0=u(n+Mdiv)+v−2Omega|
Mdiv. Then the exact rewritten inner expression is

E0\^{}2−2E0 Omega V0 V1+Omega\^{}2 V0\^{}2 Delta−Omega\^{}2 a0.

Its unique leading summand has degree 46162L+148 and
\path|coefficient2^96.| After squaring, the final subtractand
4Omega\^{}4 V0\^{}2 a0 a1 has smaller degree. Thus the exact degree is
92324L+296 and uniform leading coefficient 2\^{}192≠0. This also proves
exact total degree: these formulas give matching multivariate upper
bounds and nonzero diagonal leading terms. The same degree and leader
survive shifts to the strictly positive domain.

\subsection{Circuit identities and
accounting}\label{circuit-identities-and-accounting}

Set t\_j=B\^{}(5\^{}j), \path|v_j=B^(5^31−5^j),| with v\_31=1 and
v\_j=t\_j\^{}4 v\_\{j+1\}. With weights \path|w=(t_31,v_0,...,v_31),|

\begin{verbatim}
e(B)=t_31[z h4(w)+sum_alpha c_alpha t_31^(4−|alpha|) product_j v_j^alpha_j].
\end{verbatim}

The complete homogeneous polynomial h4 is evaluated by the ascending
recurrence h\_k←h\_k+w h\_\{k−1\}, k=1,\ldots,4. The correction is
evaluated by recursive multivariate Horner, selecting a variable and
splitting each term into one divisible by that variable and one not.
Every step is exact distributivity. The deterministic heuristic chooses
maximum incidence, then exponent mass, then a fixed permuted tie-break.
A finite search of 100 such heuristics selected8; no optimality is
claimed.

Lambda is evaluated without division as
\path|product_{j=0}^31[1+t_j(1+t_j)(1+t_j^2)](1+q)(1+q^2).| Identical
coefficient recipes become one leaf and literal recipes become literal
integers. Each of the 402 alias identities is checked exactly as an
identity in the four ports, not numerically guessed. Shared gates and
literal constant folding are allowed; every binary addition,
subtraction, or multiplication, including scaling by a numeral, is
charged. Registers, references, and fixed integer leaves are free.
Coefficient height is a separate metric, and the bitlength of y is not
hidden in a gate count.

\subsection{Reproducible checks and
limitations}\label{reproducible-checks-and-limitations}

build\_signed\_nine.py regenerates the four supplied structured/new
circuits. check\_degree.py checks cancellation-free exact leaders.
check\_signed\_nine.py compares inherited core named quantities and
exact rewrite identities modulo three primes, verifies all 402
coefficient alias identities, checks three full 58,905-term digit sums,
6,144 bound assignments including negative \path|R,C=0,H=0,86,016|
bounded finalizer assignments, and 468 explicit finalizer completeness
witnesses. These are implementation regressions, not proofs that the
final Diophantine equation has no solution. check\_actual\_index.py
reads the authentic four coefficient files, verifies hashes/bitlengths,
evaluates every coefficient bound on actual integers, and checks all 78
stored signed-magnitude coefficient payloads against their exact
recipes, then records actual-instance modular outputs.
materialize\_fixed\_leaves.py constructs those payloads. bind\_index.py
constructs the finite bound and materialized z. All programs are newly
authored or carried forward from our preceding independently authored
audit, not executed copies of received upstream programs.

\section{Fourteen unrestricted
integers}\label{fourteen-unrestricted-integers}

These are the recommended signed constructions of this package. Earlier
sixteen/seventeen-variable intermediate constructions are superseded
here. Both equations below have unsolvability provably equivalent in ZFC
to Con(ZFC). The proofs are mathematical, not newly compiled Lean or
Rocq certificates. No global minimum is asserted.

\subsection{Shared definitions and
coefficients}\label{shared-definitions-and-coefficients}

Use the exact Jones source registers and three-root finalizer in the
shared lemmas below. The source quartic has 31 positive variables;
L=5\^{}32. Every weighted coefficient is fixed by INDEX\_CONSTANTS.json,
using the whole coefficient tuple in
../coefficients/payloads/COEFFICIENT\_BINDING.json. The new index is
z=2\^{}4259933. Set

\begin{verbatim}
 Bodd=2(2z)^(L+1)=2^[1+4259934(L+1)].
\end{verbatim}

This is a fixed integer, not an unknown. Its exponent is odd, since L+1
is even. The index theorem gives \path|y=e_*(2z)<(2z)^(L+1)<Bodd.| Bodd
and y are specified by exact bounded recipes, not materialized as full
binary strings. All four native coefficient ports and z are
materialized.

\subsection{The fourteen-variable
polynomial}\label{the-fourteen-variable-polynomial}

The unknowns are exactly

\begin{verbatim}
 (epsilon,g0,h,s0,wroot,phi,i,j,z1,z2,z3,z4,residue,normroot) in Z^14.
\end{verbatim}

In the Jones source definitions set

\begin{verbatim}
 b=Bodd epsilon^2, w=wroot^2,
 g=g0+1, s=s0+1; h,i,j are unrestricted coordinates directly.
\end{verbatim}

The other registers through a0,a1,m,Ddiv,Tdiv are unchanged. Guard only
\path|x1,x2,x3=(g0,s0,phi):|

\begin{verbatim}
 d_r=3x_r+2 (r=1,2,3), d4=3m-1,
 C0=1, B_r=1+4d_r C_(r-1) z_r^2, C_r=C_(r-1)B_r,
 Cstar=C4, Z=epsilon wroot i z1 z2 z3 z4.
\end{verbatim}

Define the complete finalizer by

\begin{verbatim}
 S=Z Ddiv^5, r=Tdiv+Ddiv residue,
 T=(r^3-Ddiv^2 r-Ddiv^3)(r^2+23Ddiv^2),
 w1=Cstar, w2=Cstar(P^2+K^2+2), X=S normroot-T,
 A0=S^2 Cstar, A1=S^2 w1^2 a0, A2=S^2 w2^2 a1,
 V=X^2+A2-A0-A1, W=V^2+4X^2 A2-4A0A1.
\end{verbatim}

The equation is

\begin{verbatim}
                F14=W^2-16V^2 X^2 A2=0.
\end{verbatim}

This is fully specified by signed14.json, with 1586 binary operations:
684 additions, 26 subtractions, 876 multiplications. Its exact degree is

\begin{verbatim}
                813056*5^32+2152.
\end{verbatim}

There are 80 nonliteral fixed leaves after aliases and integer constants
are folded. Polynomial register abbreviations do not increase the 14
quantified-variable count.

\subsection{Why the square w restriction is
complete}\label{why-the-square-w-restriction-is-complete}

The source Ratio.exists\_B25core (pinned Ratio.lean.txt, necessity
section) chooses, for b=2\^{}k,

\begin{verbatim}
                w=2^(2R+1-k).
\end{verbatim}

The exponent is strictly positive: b\textless=R implies
\path|k<2^k=b<=R<2R+1.| The cofinal packing proof permits arbitrarily
large exponents k of either parity. In Master3.USys\_of\_mem lines
312--320, the selected exponent is used only to infer b\textgreater xy
and that every quartic witness coordinate is below b; the rest of the
proof uses these inequalities and b=2\^{}k, not the special formula for
that exponent. NineReduce propagates the same b unchanged. Our Bodd has
odd exponent a, and choosing epsilon=2\^{}t gives k=a+2t odd. Therefore
2R+1-k is even and the canonical w is a positive square. Choose

\begin{verbatim}
wroot=2^((2R+1-k)/2).
\end{verbatim}

All subsequent constructions in exists\_B25core and
NineRatio.ratio\_complete preserve that w. This argument uses the exact
canonical w formula; it does not assert that arbitrary Pell witnesses
may be required to be squares.

For soundness, the finalizer implies S!=0, hence epsilon,wroot,i and all
z\_r are nonzero. Consequently b\textgreater Bodd-1\textgreater y and
w\textgreater0. The guard stack gives g,s positive, phi\textgreater=0
and m\textgreater0. The h-positivity proof in the shared lemmas below
uses only these statements, not any j condition, so h\textgreater0 as
before.

\subsection{Why the j guard is
unnecessary}\label{why-the-j-guard-is-unnecessary}

The signed-H proof included below extends the product-Pell soundness
lemma to arbitrary signed H=B+jC. Its hypotheses hold here:
\path|R>0,A>1,B=2R+1>1| is \path|odd,C=B+bw+phi>=B,J=4A-5>0| and i!=0.
The quotient condition JF\textbar QF+J(H-C), together with gcd(J,F)=1,
gives J\textbar Q and F\textbar H-C. The signed-H lemma recovers
C=psi\_A(B), even when j is negative. Its key step reflects a negative
Pell index modulo 4t1, using \path|M_A^(4t1)=identity| modulo
chi\_A(t1).

The original positive product-Pell completeness theorem then supplies
fresh positive i',j' for the same A,B,C,J. Thus D F' I' is square and
F'\textbar H'-C. Since J\textbar Q is unchanged, JF'\textbar QF'+J(H'-C)
follows. Neither a0 nor m depends on i,j. Therefore the recovered signed
solution can be replaced by an ordinary positive-i,positive-j Jones
ratio solution. All other coordinates, including the restricted square
w, are unchanged. The ordinary ratio soundness theorem applies.

Conversely, the cofinal odd-exponent packing construction gives a ratio
solution with square w and positive i,j. The linear-prefix stack
supplies four nonzero guard variables, the intersective lemma supplies
residue, and the three-root finalizer supplies normroot. These arguments
prove both directions of the refutation/solution equivalence. Negation
gives

ZFC proves (there is no integer zero of F14) iff Con(ZFC).

\subsection{A lower-degree equation in fifteen
integers}\label{a-lower-degree-equation-in-fifteen-integers}

Keep b=y+1+epsilon, with epsilon guarded rather than squared, but still
set w=wroot\^{}2 and release j. Guard the four coordinates
(epsilon,g0,s0,phi) and m, using five z variables. Set Z=wroot i z1 z2
z3 z4 z5. The remaining definitions are unchanged. This gives
signed15\_linear.json, with 1594 operations (688 additions, 26
subtractions, 880 multiplications) and degree

\begin{verbatim}
                406528*5^32+2528.
\end{verbatim}

Completeness chooses a sufficiently large odd k in the free-base packing
proof, then \path|epsilon=2^k-y-1>=0.| The canonical w is again a
square. Soundness needs only b\textgreater y and w\textgreater0; it does
not require an independently asserted parity of k. Using h,i,j directly
rather than h0+1,i0+1,j0+1 is an affine bijection on unrestricted
integers and saves three gates without changing degree. This form has
roughly half the degree of F14, at the cost of one additional integer
variable and eight operations.

\subsection{Exact degree certificates}\label{exact-degree-certificates}

For F14 the named degrees are:

m:4442L+14; Cstar:4442L+58; Ddiv:13750L+29; S:68750L+152; w2:8242L+68;
A2:203264L+538.

A2 strictly dominates X\^{}2,A0,A1, and the second term in the final
subtraction has smaller degree than A2\^{}4. On the diagonal its leading
coefficient is \path|2^832*12^120*Bodd^(406528L+8),| which is nonzero.
For F15linear the corresponding output coefficient is
2\^{}832*12\^{}248, and \path|deg(A2)=101632L+632.| These give the
stated exact degrees for the actual coefficients, not merely generic
degree bounds. The independent modular checks supplement these exact
mathematical leading-term arguments.

\subsection{Evidence and limitations}\label{evidence-and-limitations}

check\_signed\_forms.py regenerates each circuit from inert source data,
checks source-register equality on signed inputs, and independently
computes the finalizer as an8-by-8 multiplication-matrix determinant.
Finite guard, intersectivity and divisibility tests are included. The
universal proofs, not the finite tests, establish equivalence. No actual
ZFC refutation or halting witness is searched for, and no consistency
claim is made.

Further reductions would require proving a further sign restriction is
redundant or a further canonical witness lies in a polynomially
parameterized subfamily. Neither arbitrary square substitution nor
finite testing would justify such a step.

\subsection{Complete shared Jones source
registers}\label{complete-shared-jones-source-registers}

\begin{verbatim}
B=b^5, q=B^L, ell=sum_{r=1}^{31}B^(5^r),
lambda=sum_{r=0}^{4L-1}B^r, e=e_*(B), theta=B-2z,
N=q^16, Gcap=32b B^(5^31),
\end{verbatim}

\begin{verbatim}
R=[g+e q^3+ell q^5+
   (2(e-z lambda)(1+B+g)^4+lambda B+lambda B q^4)q^7](N^2-N)
  +[q^3-b ell+ell+theta lambda q^3+(B-2)q^8](N^2-1).
\end{verbatim}

\begin{verbatim}
U=N^2w, Y=N^2s, M=RY, P=2M^2U,
K=R+1+h(P-1), A=M(U+1), B1=2R+1, W1=bw,
C=B1+W1+phi, J=4A-5, Q=3W1 C-2(W1^2-1),
D=(A^2-1)C^2+1, E=iJDC^2, F=(A^2-1)E^2+1,
Gp=A+F(CD+1-A), H=B1+jC, I=(Gp^2-1)H^2+1,
a0=(P^2-1)K^2+1, a1=DFI,
m=K^2 Gcap-4(C-KY)^2 Gcap-8gK^2,
Ddiv=JF, Tdiv=QF+J(H-C).
\end{verbatim}

\subsection{Shared algebraic lemmas}\label{shared-algebraic-lemmas}

\subsection{Linear-prefix square
stack}\label{linear-prefix-square-stack}

For any integer assignment, every B\_r and C\_r is 1mod4, hence nonzero,
and \path|gcd(B_r,C_(r-1))=1.| Thus all factors B\_r are pairwise
coprime. If all z\_r!=0 and Cstar is an integer square, each
\textbar B\_r\textbar{} is a square. As B\_r=1mod4, a negative B\_r
cannot be a signed negative square, so every B\_r is a positive square.
All prefix C\_r are consequently positive. If any d\_r\textless0, its
factor would be \textless=-3, a contradiction. No d\_r vanishes because
d\_r=2mod3. Hence d\_r\textgreater0: all guarded x\_j\textgreater=0 and
m\textgreater0.

Conversely, suppose these inequalities hold. Inductively C\_(r-1) is a
positive square. The positive integer d\_r is 2mod3 and nonsquare, so
d\_r C\_(r-1) is nonsquare. Choose a nontrivial Pell solution
u\^{}2-d\_r C\_(r-1) v\^{}2=1 with v even (a suitable positive power of
a Pell unit has even second coordinate). Put z\_r=v/2!=0. Then
B\_r=u\^{}2 and C\_r is a positive square. This constructs the guard
list.

When S!=0, all z\_r!=0. Each d\_r is nonzero and each prefix nonzero, so
B\_r!=1; since B\_r=1mod4, \textbar B\_r\textbar\textgreater=3.
Therefore \textbar Cstar\textbar\textgreater=3\^{}r\textgreater1, where
r is the number of guards. This observation permits the smaller signed
separation weight w1=Cstar.

\subsection{Divisibility and
nonvanishing}\label{divisibility-and-nonvanishing}

For every integer base assignment Ddiv!=0. Indeed J=4A-5!=0. If F=0,
then (A\^{}2-1)E\^{}2=-1 forces A=0,E=+/-1, but E=iJDC\^{}2 is then
divisible by 5, impossible.

The polynomial \path|Pi(t)=(t^3-t-1)(t^2+23)| has no rational root and
has a root modulo every positive integer. The cubic is monic and
neither1 nor-1 is a root; the quadratic has no real root. For odd primes
p!=23, either the cubic has a simple root, or it is an irreducible cubic
whose Frobenius three-cycle fixes a square root of its discriminant -23;
then the quadratic has a simple root. Lift by Hensel. At23 the cubic's
root3 is simple. At2 the quadratic lifts because -23=1mod8. Apply CRT.

For nonzero D and arbitrary integer T0,Z,

D\textbar T0 and Z!=0 iff exists v: ZD\^{}5 \textbar{}
\path|(r^3-D^2r-D^3)(r^2+23D^2),| r=T0+Dv.

The right-hand polynomial never vanishes, since it equals D\^{}5
Pi(T0/D+v). If D\textbar T0, intersectivity supplies v modulo
\textbar Z\textbar. Conversely Z=0 would force a zero. If p-adic
valuations satisfy \path|v_p(T0)=beta<alpha=v_p(D),| then v\_p(r)=beta,
the cubic and quadratic factors have valuations3beta and 2beta
respectively, so their product has valuation5beta\textless5alpha,
contradicting divisibility byD\^{}5. Apply with D=Ddiv,T0=Tdiv.

\subsection{Three-root finalizer with signed small
weights}\label{three-root-finalizer-with-signed-small-weights}

The polynomial F is identically the product of the eight factors X-S(e0
sqrt(Cstar)+e1 w1 sqrt(a0)+e2 w2 sqrt(a1)), e\_r in\{-1,1\}. Complex
square roots are allowed in this formal identity. T is never zero. If
S=0, then F=T\^{}8!=0. Hence every zero has S!=0.

On this locus, \path!|w1|=|Cstar|>sqrt(|Cstar|).! Also, for arbitrary
real P,K, \path!sqrt(|(P^2-1)K^2+1|)<=P^2+K^2+1.! Therefore
\textbar w2\textbar\textgreater sqrt(\textbar Cstar\textbar)+\textbar w1\textbar sqrt(\textbar a0\textbar).
Distinct sign patterns of the nonzero roots have distinct sums: the
largest differing root has magnitude at least1 (its square is a nonzero
integer), and its weighted contribution dominates all earlier ones. If
F=0, one sum equals normroot-T/S, a rational number. Every Galois
automorphism of the field generated by these roots changes their signs;
the separation forces it to fix each root. Each root is rational, and
then integer since its square is integer. Thus Cstar,a0,a1 are squares
and S\textbar T. This argument does not assume independent square
classes and works for initially negative radicands. Conversely those
square and divisibility conditions allow
\path|normroot=T/S+r0+w1r1+w2r2.|

\subsection{Removing the h and i sign
restrictions}\label{removing-the-h-and-i-sign-restrictions}

The guards and square w give g,s,w positive,
phi\textgreater=0,m\textgreater0, b\textgreater y and i!=0. Replacing i
by \textbar i\textbar{} changes none of a0,a1,Ddiv,Tdiv,m: i occurs
through E and thereafter only E\^{}2. Hence i can be taken positive.

To show h\textgreater0, m\textgreater0 implies K!=0,
g\textless Gcap/8=4b B\^{}(5\^{}31), and 4(C-KY)\^{}2\textless K\^{}2.
Since b\textgreater y, B\textgreater2z and B\textgreater=32. With
t=5\^{}31, 1+B+g \textless{} B\^{}(t+2)\textless=q and b ell\textless q.
These bounds make both brackets defining R positive (the positive lambda
B q\^{}4 absorbs 2z \path|lambda(1+B+g)^4).| Thus
R\textgreater=N\^{}2-N\textgreater=2. Consequently
C\textgreater0,Y,U\textgreater=1, P\textgreater=2R\^{}2\textgreater R+2.
The strict approximation forces K\textgreater0. If h\textless=0, then
1\textless=K\textless=R+1\textless P. K=1 would imply P-1\textbar R,
impossible. For2\textless=K\textless P,
\path|(PK-1)^2<(P^2-1)K^2+1<(PK)^2,| contradicting a0 square. Hence
h\textgreater=1.

\subsection{Removing the Jones j-positivity guard by signed index
reflection}\label{removing-the-jones-j-positivity-guard-by-signed-index-reflection}

Let A,B,C,J be positive integers with A\textgreater1, B\textgreater1,
B\textless=C, J\textgreater0, and assume 2B\textless=C or B is odd. Let
i be any nonzero integer and j any integer. Define

\begin{verbatim}
D=(A^2-1)C^2+1, E=i J D C^2, F=(A^2-1)E^2+1,
G=A+F(CD+1-A), H=B+jC, I=(G^2-1)H^2+1.
\end{verbatim}

Suppose DFI is an integer square and F divides H-C. Then C=psi\_A(B),
where psi\_A is the Pell sine sequence with recurrence psi\_A(n+2)=2A
psi\_A(n+1)-psi\_A(n), psi\_A(0)=0, psi\_A(1)=1.

\subsection{Proof}\label{proof}

Replacing i by \textbar i\textbar{} leaves D,F,G,H,I unchanged. Thus E
may be assumed positive. The standard product-Pell argument gives D,F,I
positive, G\textgreater1, and pairwise coprime D,F,I. Indeed F=1 modulo
D, G=1 modulo D, I=1 modulo D, G=A modulo F, and H=C modulo F imply I=D
modulo F. Hence all three factors are positive squares.

The Pell parametrization supplies nonnegative indices t0,t1,t2 with
C=psi\_A(t0), E=psi\_A(t1), \textbar H\textbar=psi\_G(t2). As in the
pinned ProductPell proof, C\^{}2\textbar E implies C\textbar t1; also
t1\textgreater=2, 1\textless=t0\textless=t1, and F=chi\_A(t1)\^{}2. Set
k=chi\_A(t1)\textgreater0.

If H\textgreater=0, the original proof applies verbatim: H=psi\_G(t2),
psi\_A(t2)=C modulo k, and t2=B modulo C. The usual Pell
index-congruence lemma gives B=t0 or B=-t0 modulo C, and the
productPell\_index lemma gives C=psi\_A(B).

If H\textless0, then psi\_A(t2)=-C modulo k. Let
\path|M=[[2A,-1],[1,0]].| The determinant of M is 1 and the trace of
M\^{}t1 is 2chi\_A(t1)=0 modulo k. Cayley-Hamilton therefore gives
M\^{}(2t1)=-Id modulo k, and M\^{}(4t1)=Id modulo k. Choose a positive
integer Q large enough that n=4t1 Q-t2\textgreater=0. The inverse-matrix
identity for the Pell sine coordinate gives

\begin{verbatim}
psi_A(n)=-psi_A(t2)=C modulo k.
\end{verbatim}

Since G=1 modulo C, psi\_G(t2)=t2 modulo C. Thus H=-psi\_G(t2) implies
-t2=B modulo C. As C\textbar t1, the reflected index satisfies n=-t2=B
modulo C. Apply the same ordinary index-congruence lemma now to n in
place of t2. Its conclusions are n=t0 or n=-t0 modulo C, so again B=t0
or B=-t0 modulo C, and the unchanged productPell\_index conclusion
follows. This proof does not require j\textgreater=0 or H\textgreater=0.
The case H=0 causes no issue for the nonnegative branch; alternatively
F\textgreater C and F\textbar H-C exclude it.

\subsection{Application to the Jones ratio
conditions}\label{application-to-the-jones-ratio-conditions}

Retain the guards on g,s,w,phi and the positivity consequence for h.
They imply \path|R>0,A>1,C>=B=2R+1>1,| with B odd. Keep i!=0 but drop
the guard j0\textgreater=0: put j=j0+1 with j0 an arbitrary integer.
From JF\textbar QF+J(H-C) and gcd(J,F)=1, derive J\textbar Q and
F\textbar H-C. The signed-H lemma yields C=psi\_A(B). The standard
constructive product-Pell theorem then supplies fresh positive i',j' for
these same A,B,C,J. Their F',H' satisfy F'\textbar H'-C. Together with
J\textbar Q this restores JF'\textbar QF'+J(H'-C), and D F' I' is
square. The radicand a0 and margin m do not depend on i,j, so they
remain unchanged. Thus any signed-j ratio solution can be replaced by a
genuine positive-j ratio solution. Completeness is immediate because the
original positive witnesses are still allowed.

This removes one guard witness from the coprime Pell stack. It uses
ordinary Pell theory and matrix arithmetic, not a newly compiled formal
theorem. The exact source proof inspected is
Computability/HilbertTenthProblem/Lean/Diophantine/Paper1982/ProductPell.lean,
psi\_of\_ProductPellConds and productPell\_index, at commit
\path|8c6eec9d19ef80538c8e61ff6e0180040583446a|.

\section{Verification and
reproducibility}\label{verification-and-reproducibility}

The archive records a pinned source ledger and exact arithmetic
circuits. Its root verifier first checks immutable file hashes, then
runs the new checks in a temporary copy. Source programs received from
the repository are preserved as inert text or data; no received program
is executed by the release verifier. The archive includes the symbolic
identities and full proof text rather than presenting random evaluation
as proof.

The finite checks include independent reconstruction of the coefficient
payloads and controller hypotheses; exact live-gate counts; the quartic
register identity; formal multiplication of the radical norm factors;
formal-Laurent degree checks for all positive variants; independent
three-root determinant evaluation; and signed Pell matrix-reflection
identities. Complete machine-readable receipts and a check plan are
included.

Some bounded signed-Pell searches have no admissible hits. These are not
reported as evidence for a universal theorem; the mathematical
reflection proof establishes that result. Actual fixed-coefficient
evaluations and uniform nonzero leading-term proofs are distinguished
from diagnostic modular specializations.

A positive computational result here means that the stated finite check
passed. It is not a successful search for a solution of the Con(ZFC)
equation. There is no such search or witness claimed in this report.

\section{Further research questions}\label{further-research-questions}

\begin{enumerate}
\def\labelenumi{\arabic{enumi}.}
\tightlist
\item
  Can the remaining signed Pell guards be reduced without restricting
  the Jones ratio witnesses to a family lacking a completeness theorem?
\item
  Can a general formal theorem expose the cofinal choice of Jones radix,
  so that square-scaled substitutions no longer require inspection of
  the existing completeness proof?
\item
  Can the carrier translations be optimized systematically, with a
  symbolic valuation certificate tracking every radical branch and its
  cancellation?
\item
  Can another quartic register be eliminated by an identity modulo
  retained equations, rather than by direct substitution that raises
  degree?
\item
  Can the new exact coefficient layout and actual-rank controller be
  formalized once as a reusable compiler interface?
\item
  Can coefficient expressions be reduced further without concealing
  their materialization cost or changing the meaning of the operation
  metric?
\item
  Can an explicit instantiation of a smaller unrestricted-integer
  universal representation improve this archive's count while retaining
  a complete coefficient recipe and a manageable proof?
\end{enumerate}

\section{Sources attribution and
limitations}\label{sources-attribution-and-limitations}

The primary source is the pinned ProveIt catalogue and its packages
24--30:

\url{https://github.com/VladimirReshetnikov/ProveIt/tree/8c6eec9d19ef80538c8e61ff6e0180040583446a/SetTheory/Cardinals/docs/reports/hilbert-tenth-problem/con-zfc-diophantine-forms}

The source ledgers identify exact files and Git blob hashes. The
mathematical contributions of the earlier packages, Jones's reduction,
the K215 construction, and the compiler ancestry remain credited in the
retained notices. This report and its newly authored checks are
OpenAI-assisted research prepared for Vladimir Reshetnikov's ProveIt
project. Independent audits described here are separate checks within
that collaboration, not external human endorsements.

New material follows ProveIt's MIT No Attribution terms. Inherited
third-party material retains its own licence; the Wade licence and K215
attribution notice are included. Historical source notices may name
paths in their original packages. The active release paths and
coefficient bindings are defined by the root README and binding
manifests, not by those historical references.

Sun's 2026 preprint states in Theorem 1.1 a representation of every
recursively enumerable subset of the natural numbers using ten
unrestricted integer witnesses. This prevents any interpretation of our
count as a global fewest-variable record. Our contribution is an
explicitly instantiated polynomial with a fixed coefficient recipe and a
detailed dependency chain; we have not instantiated Sun's general
construction for the particular ZFC-refutation predicate. See
\url{https://arxiv.org/html/2609.01594v2} and the archive's scope note
for the exact distinction. We have not independently audited the
complete proof of that external preprint.

There is no global minimality theorem, proof of Con(ZFC), new
kernel-checked end-to-end result, or claim that every large fixed
numeral is materialized.

\end{document}
